\chapter{Autonomous Movement}
\label{ch:movement}

During the autonomous period, nobody is steering. The robot can only follow a
list of instructions such as ``drive 1000~mm forward, then turn 90 degrees
right''. This chapter explains how the program turns those instructions into
motor rotations, how we calibrate the result, and how the final autonomous
routine uses them.

\section{Encoders and Wheel Degrees}
\label{sec:encoders}

Each V5 motor contains an \emph{encoder}, a sensor that counts how far the motor
shaft has turned, in degrees. The motor can be told ``turn exactly 1128
degrees'' and it will stop itself once the encoder reaches that number. The
robot has no sensor that measures distance along the floor directly, so every
autonomous movement must be converted into a number of shaft degrees first.

\section{The Original Approach}
\label{sec:original}

The program originally used the SDK's built-in \code{DriveTrain} class, which
does this conversion internally. It was constructed as follows (from commit
\commit{4f76b82}, since removed):

\begin{lstlisting}[caption={Original drivetrain construction, \code{src/main.py} at \commit{4f76b82}.},label=lst:olddt,numbers=none]
dt = DriveTrain(left_motor, right_motor, 300, 320, 320, MM)
\end{lstlisting}

The three numbers are the wheel travel, track width and wheelbase in
millimetres. However, these are the SDK's example values, not measurements of
our robot. Our wheels are 4-inch (101.6~mm) wheels, which travel about 319~mm
per revolution, not 300~mm. Since every movement depended on these numbers,
every autonomous distance and angle was wrong. We replaced \code{DriveTrain}'s
movement functions with our own maths (commit \commit{eb032a0}), so that every
step of the conversion is visible and can be checked.

\section{Driving Straight}
\label{sec:drive}

When a wheel turns once, the robot moves forward by the wheel's circumference,
$\pi D$, where $D$ is the wheel diameter. To move a distance $d$, the wheel must
therefore turn $d / (\pi D)$ revolutions, which is $360\,d/(\pi D)$ degrees.
Including the external gear ratio $G$ and a correction factor $S_d$, the number
of shaft degrees is
\begin{equation}
  \theta_{\text{drive}} = \frac{d}{\pi D} \times 360 \times G \times S_d .
  \label{eq:drive}
\end{equation}

$G$ is the number of motor revolutions per wheel revolution. Our wheels are
mounted directly on the motors, so $G = 1$. $S_d$ is \code{DRIVE_SCALE}, a
multiplier that absorbs small errors such as tyre squash or a slightly wrong
diameter. It starts at 1.00. Equation \eref{eq:drive} is implemented directly:

\lstinputlisting[firstline=9,firstnumber=9,lastline=11,caption={\code{wheel_degrees_for_distance()} in \code{src/autonomousroutine.py}, lines 9--11.},label=lst:drivedeg]{\src{autonomousroutine.py}}

\textbf{Worked example:} To drive $d = 1000$~mm with $D = 101.6$~mm,
\[
  \pi D = \pi \times 101.6 \approx 319.2~\text{mm},
  \qquad
  \theta_{\text{drive}} = \frac{1000}{319.2} \times 360 \approx 1128^\circ .
\]
So each wheel turns just over three full revolutions. The same function with
the SDK's default value of 300~mm would give $1200^\circ$, and the robot would
drive about 1064~mm instead of 1000~mm, an error of over 6\%.

\section{Turning in Place}
\label{sec:turn}

To turn on the spot, the two wheels spin in opposite directions. The robot then
rotates around the point halfway between its wheels, and each wheel travels
along a circle whose diameter is the \emph{track width} $W$, the distance between
the left and right wheels (\fref{fig:turn}).

\begin{figure}[h]
  \centering
  \begin{tikzpicture}[>={Stealth}, scale=0.9]
    \draw[dashed] (0,0) circle (2);
    \draw[thick, fill=codebg] (-1.3,-1.6) rectangle (1.3,1.6);
    \fill (0,0) circle (2pt) node[below right] {\small centre};
    \draw[fill=gray!60] (-2.25,-0.6) rectangle (-1.75,0.6);
    \draw[fill=gray!60] (1.75,-0.6) rectangle (2.25,0.6);
    \draw[<->] (-2,-2.4) -- node[below] {\small track width $W$} (2,-2.4);
    \draw[->, very thick, codekw] (-2.7,-0.6) -- (-2.7,0.8);
    \node[left, font=\small, align=right] at (-2.8,0.1) {left wheel\\forward};
    \draw[->, very thick, codest] (2.7,0.6) -- (2.7,-0.8);
    \node[right, font=\small, align=left] at (2.8,-0.1) {right wheel\\backward};
    \draw[->, thick] (110:2.5) arc (110:70:2.5);
    \node[font=\small] at (0,2.85) {robot turns right by $\alpha$};
  \end{tikzpicture}
  \caption{A right point turn, seen from above. Each wheel follows the dashed circle of diameter $W$ (self-made).}
  \label{fig:turn}
\end{figure}

One full spin moves each wheel along the whole circle, $\pi W$. A turn of
$\alpha$ degrees is the fraction $\alpha/360$ of a full spin, so each wheel
travels an arc of length
\begin{equation}
  s = \pi W \times \frac{\alpha}{360} .
  \label{eq:arc}
\end{equation}

Note: this is not doubled. It can look as if both wheels together must cover
$2\pi W$, but each wheel only covers its own half of the motion; the two sides
split the circle by moving in opposite directions. \code{AGENTS.md} warns
about this because it has been mistaken for an error before.

Converting the arc $s$ into shaft degrees in the same way as \eref{eq:drive},
with the turn correction factor $S_t$ (\code{TURN_SCALE}),
\begin{equation}
  \theta_{\text{turn}} = \frac{\pi W \alpha / 360}{\pi D} \times 360 \times G \times S_t
                       = \frac{W}{D}\,\alpha\, G\, S_t .
  \label{eq:turn}
\end{equation}

The $\pi$ and the 360 cancel, which leaves a neat result: the wheel turns
$W/D$ times as far as the robot does.

\lstinputlisting[firstline=13,firstnumber=13,lastline=16,caption={\code{wheel_degrees_for_turn()} in \code{src/autonomousroutine.py}, lines 13--16.},label=lst:turndeg]{\src{autonomousroutine.py}}

\textbf{Worked example:} For a 90° right turn, with $W = 320$~mm,
\[
  \theta_{\text{turn}} = \frac{320}{101.6} \times 90 \approx 3.15 \times 90 \approx 283^\circ .
\]
The left wheel turns $283^\circ$ forward and the right wheel $283^\circ$
backward. \code{turn()} does this by passing the same number to both sides with
opposite signs:

\lstinputlisting[firstline=39,firstnumber=39,lastline=47,caption={\code{drive()} and \code{turn()} in \code{src/autonomousroutine.py}, lines 39--47.},label=lst:driveturn]{\src{autonomousroutine.py}}

The right motor is declared \emph{reversed} in \code{config.py} (it faces the
other way on the robot), so ``forward'' means robot-forward on both sides. That
is why a negative number on the right side is what makes the robot turn.

\section{Stopping Accurately: HOLD versus BRAKE}
\label{sec:hold}

A V5 motor has three ways of stopping:

\begin{itemize}
  \item \textbf{COAST} cuts the power and lets the motor spin freely until
    friction stops it.
  \item \textbf{BRAKE} shorts the motor so it resists turning, which stops it
    quickly, but it stops resisting once it has stopped.
  \item \textbf{HOLD} actively pushes the motor back to its target position, even
    after it has arrived.
\end{itemize}

The original program used \code{BRAKE} during autonomous. The problem is that a
moving robot has momentum, so it slides slightly past its target. With
\code{BRAKE}, this overshoot is kept, and it adds up across every movement in the
routine. With \code{HOLD}, the motor notices that it has passed the target and
pulls back. \code{AGENTS.md} records that switching to \code{HOLD} removed most
of the original turn error. \code{move_wheels()} therefore sets \code{HOLD} before
every autonomous movement (lines 28--29 of \code{autonomousroutine.py}).

Driver control still uses \code{BRAKE}, because a driver expects the robot to
stop and stay where it stops, not pull back towards an old target.

\section{Timeouts}
\label{sec:timeouts}

Every movement has a \emph{timeout}, a maximum time after which the motor gives
up. This prevents the robot from pushing against an obstacle forever. The
original program used a flat 5-second timeout for every move. However, a long
drive can need more than 5 seconds, and the timeout then silently stops it
halfway with no error.

\textbf{Worked example:} A green-cartridge motor turns at up to 200~rpm, which is
$200 \times 360 / 60 = 1200$ degrees per second (\code{MOTOR_MAX_DPS}). At 50\%
speed, it turns 600 degrees per second. \code{drive(5000)} needs
$5000 / 319.2 \times 360 \approx 5639^\circ$, which takes about 9.4~s. In
5~s, the motor only reaches $5 \times 600 = 3000^\circ$, which is about 2660~mm,
or 53\% of the requested distance.

The timeout is now calculated from the movement itself:
\begin{equation}
  t_{\text{timeout}} = 2 \times \frac{\theta}{\omega} + 1~\text{s},
  \label{eq:timeout}
\end{equation}
where $\omega$ is the expected speed in degrees per second. The factor of 2
allows for the motor being slower than its theoretical speed, and the extra
second covers the time spent speeding up on short moves.

\lstinputlisting[firstline=18,firstnumber=18,lastline=21,caption={\code{move_timeout_ms()} in \code{src/autonomousroutine.py}, lines 18--21.},label=lst:timeout]{\src{autonomousroutine.py}}

For \code{drive(5000)} at 50\%, this gives $2 \times 9.4 + 1 = 19.8$~s, which is
long enough to complete the move but still stops a robot that is genuinely
stuck.

\section{Calibration}
\label{sec:calibration}

Equations \eref{eq:drive} and \eref{eq:turn} are only as accurate as the
measurements of $D$ and $W$. In practice, turns are also affected by
\emph{wheel scrub}: during a point turn the tyres slide sideways slightly, so the
encoders reach their target while the robot has turned less than intended.
Scrub cannot be predicted from the geometry, so instead of measuring it, we
absorb it into \code{TURN_SCALE}.

If we ask for a turn of $\alpha_{\text{asked}}$ and the robot actually turns
$\alpha_{\text{measured}}$, the corrected scale is
\begin{equation}
  S_{\text{new}} = S_{\text{old}} \times \frac{\alpha_{\text{asked}}}{\alpha_{\text{measured}}} .
  \label{eq:scale}
\end{equation}

\textbf{Worked example:} We ask for a full 360° spin with $S_t = 1.00$ and the
robot stops at 330°. Then $S_{\text{new}} = 1.00 \times 360/330 \approx 1.091$, so
every future turn is made about 9\% larger to compensate.

Measuring angles by eye is slow, so we wrote \code{turn_tuner()} (commit
\commit{eb032a0}) to do the calibration on the robot itself, using the
controller. Button A spins the robot one or more full turns and B spins it back.
The up and down arrows change \code{TURN_SCALE} by 0.010, and R1 and R2 change
it by 0.001 for fine adjustment. The driver keeps adjusting until the robot ends
a full spin facing exactly the way it started, then copies the value shown on
the screen into \code{config.py}. Button Y changes the test speed, because scrub
depends on speed; the scale should be tuned at the speed the routine actually
uses.

\section{Squaring to the Wall}
\label{sec:wallsquare}

Even with perfect calibration, an autonomous routine is only as good as its
starting position. If the robot starts angled by 3°, it will end up about
52~mm off course after driving 1~m. Human placement at the start of a match is
not precise enough to avoid this.

\code{square_to_wall()} solves this by using the field wall as a reference. The
robot reverses slowly into the wall. When one side touches first, that side
stops while the other keeps going, which rotates the robot until its back is
flat against the wall.

\lstinputlisting[firstline=168,firstnumber=168,lastline=184,caption={\code{square_to_wall()} in \code{src/autonomousroutine.py}, lines 168--184.},label=lst:wall]{\src{autonomousroutine.py}}

The robot knows it has reached the wall when both motors have \emph{stalled},
meaning their speed has dropped below 5\% even though they are still powered.
Three timings make this reliable:

\begin{itemize}
  \item \code{WALL_SQUARE_MIN_MS} (400~ms): stall checks are skipped at first,
    because a motor that has only just started is also moving slowly and would
    look stalled.
  \item \code{WALL_SQUARE_PUSH_MS} (300~ms): after both sides stall, the robot
    keeps pushing briefly, so that both sides are pressed flat rather than just
    touching.
  \item \code{WALL_SQUARE_MAX_MS} (2500~ms): if no wall is found, the robot
    gives up rather than reversing forever.
\end{itemize}

\section{The Autonomous Routine}
\label{sec:routine}

The final routine follows the path in \fref{fig:path}. The robot starts at S
with its back towards the wall, facing A.

\begin{figure}[h]
  \centering
  \begin{tikzpicture}[>={Stealth}, scale=1]
    \fill[gray!35] (-1.0,-0.55) rectangle (1.0,-0.35);
    \node[font=\small] at (0,-0.85) {wall};
    \node[circle, draw, fill=white, inner sep=2pt] (S) at (0,0) {S};
    \node[circle, draw, fill=white, inner sep=2pt] (A) at (0,3.5) {A};
    \node[circle, draw, fill=white, inner sep=2pt] (B) at (3.5,3.5) {B};
    \draw[->, thick] (S) -- node[left, font=\small] {1000 mm} (A);
    \draw[->, thick] (A) -- node[above, font=\small] {1000 mm} (B);
    \draw[->] (0.45,3.05) arc (-90:0:0.45);
    \node[font=\small] at (1.3,2.85) {turn 90° right};
  \end{tikzpicture}
  \caption{The S--A--B autonomous path (self-made).}
  \label{fig:path}
\end{figure}

\lstinputlisting[firstline=186,firstnumber=186,lastline=208,caption={Arm helpers and \code{autonomous()} in \code{src/autonomousroutine.py}, lines 186--208.},label=lst:auto]{\src{autonomousroutine.py}}

The steps are:

\begin{enumerate}
  \item Square to the wall, so the robot starts straight (\sref{sec:wallsquare}).
  \item Lift the arm 90 motor degrees and \code{HOLD} it, so it does not drag on
    the floor while driving.
  \item Drive 1000~mm to A, turn 90° right, and drive 1000~mm to B.
  \item Release the arm by switching it to \code{COAST}, so it drops under its
    own weight.
\end{enumerate}

A 300~ms \code{settle()} between moves lets the robot stop rocking, so that the
momentum of one move does not carry into the next.

\section{Summary}

\begin{enumerate}
  \item Distances and angles are converted into shaft degrees with
    \eref{eq:drive} and \eref{eq:turn}, using measured wheel and track sizes
    rather than the SDK's defaults.
  \item \code{HOLD} and proportional timeouts make each move end where it
    should (\sref{sec:hold}, \sref{sec:timeouts}).
  \item \code{TURN_SCALE} absorbs wheel scrub and is tuned on the robot with
    \code{turn_tuner()} (\sref{sec:calibration}).
  \item Squaring to the wall gives a repeatable starting position for the
    S--A--B routine (\sref{sec:wallsquare}, \sref{sec:routine}).
\end{enumerate}

We now explain the driver control half of the program.
